\section{Introduction}
\label{sec:intro}
% -------------------------------------------------------------

Vision-language model (VLM) agents capable of operating a web browser
are moving from research prototypes into production deployment
% CITE: copy any line below into Semantic Scholar / Google Scholar / arXiv to find the paper.
%   - "OpenAI Operator browser-using agent 2025 launch"
%   - "Anthropic Claude computer use 2024 system card"
%   - "Google Gemini agent Project Mariner browser 2024 2025"
%   - "survey large language model web-browsing agents 2024 2025"
%   - "multimodal LLM agent benchmark survey browser navigation"
\citep{OpenAI2025Operator,anthropic2024computeruse,tamkin2024clioprivacypreservinginsightsrealworld,andreux2025surferhmeetsholo1costefficient,zhang2025surveylargelanguagemodel,nguyen-etal-2025-gui}. To complete the tasks they are
given (applying for jobs, purchasing goods, paying bills, replying to
email), these agents are routinely provisioned with the user's full
personally identifiable information (PII): legal name, home address,
payment card details, government identifiers, account credentials, and
API keys. The agent stores this profile in its context and draws on
it whenever a form field must be filled or a credential must be
submitted. The result is a system that combines broad web access with
privileged read access to high-value personal data, a configuration
that realises a \emph{semantic} variant of the classic
\emph{confused-deputy} vulnerability
% CITE: copy any line below into Semantic Scholar / Google Scholar / arXiv to find the paper.
%   - "Hardy 1988 confused deputy or why capabilities might have been invented"
%   - "Norm Hardy confused deputy ACM SIGOPS Operating Systems Review 1988"
%   - "confused deputy LLM agent tool use indirect prompt injection"
%   - "agent confused deputy untrusted input privileged context 2024 2025"
\citep{Hardy1988TheCD,Ferrag2025FromPI,Triedman2025MultiAgentSE,fu2024impromptertrickingllmagents,kim2025promptflowintegrityprevent,ji2026tamingvariousprivilegeescalation}: privileged user context and unprivileged web input share one decision loop, but exploitation is semantic---social engineering bypasses the model's behavioural safety constraints rather than escalating a privilege boundary. Untrusted content encountered while browsing can manipulate the agent into forwarding that data to an adversary. The web already hosts extensive phishing and manipulation infrastructure (look-alike domains, dark patterns, fake trust signals, urgency timers, authority impersonation, conversational escalation) built to harvest PII from humans
% CITE: copy any line below into Semantic Scholar / Google Scholar / arXiv to find the paper.
%   - "APWG Anti-Phishing Working Group quarterly phishing activity trends report 2024 2025"
%   - "Mathur dark patterns shopping websites taxonomy CHI 2019"
%   - "Gray dark patterns side of UX design CHI 2018"
%   - "Brignull deceptive design patterns taxonomy"
%   - "Bösch dark strategies online services privacy"
%   - "Verizon Data Breach Investigations Report 2024 social engineering"
\citep{Hamadouche_Boudraa_Gasmi_2024,10.1145/3359183,10.1145/3173574.3174108,deceptivedesignpatterns,Bsch2016TalesFT}. Unlike a human, who can pause or close the tab, an agent is optimised to finish its task and will follow links, fill fields, and submit forms. Capability benchmarks confirm that frontier agents are now reliable enough to do this end-to-end: REAL \citep{garg2026real} reports 41\% success on multi-turn state-changing web tasks---the same competence that makes agents useful makes them a reliable target for data exfiltration. We find that even when an independent LLM judge confirms the agent's reasoning has flagged a site as suspicious, the agent still hands over critical PII at a high rate---a detection--action gap between what the agent \emph{says} and what it \emph{does}, robust across all four model families (\S\ref{sec:results-f1}).

\paragraph{The measurement gap.} Despite a growing literature on the
adversarial robustness of web agents, existing evaluation frameworks
gate success on proxy outcomes (click compliance, indirect prompt
injection at tool-use APIs, single-click redirects, or incidental
over-sharing on benign sites) rather than on the terminal field-level
submission of sensitive data to an attacker-controlled endpoint
(Table~\ref{tab:related-matrix}, \S\ref{sec:related})
\citep{cuvin2026decepticon,ersoy2026trickyarena,chen2026obviousinvisiblethreat,debenedetti2024agentdogo,korgul2025trap,zharmagambetov2026agentdam,xinyi2026webtrappark}.
The terminal submission is the
data-loss event a deployment review cares about: early-trajectory
anomalies are observable to low-cost runtime guardrails, but a
completed form-level POST payload is a terminal, unrecoverable
leakage event that only the model's own refusal can prevent. To our knowledge,
no published benchmark combines (i) attacker-controlled websites
that explicitly social-engineer the agent, (ii) a populated PII
profile spanning critical, high, and medium sensitivity tiers,
(iii) field-level exfiltration as the primary outcome, and (iv)
axis-controlled siblings that isolate which design factors causally
drive leakage.

\paragraph{This work.} We introduce \textsc{Scammer4U}, a benchmark
of 91 hand-crafted attacker-controlled environments plus 10
benign-twin baselines, designed to answer a single operational
question: \emph{when an autonomous web agent encounters a
social-engineering attack while browsing, will it hand over the
user's PII?} Environments cover eight attack vectors and span 16
site categories (full enumeration in \S\ref{sec:vectors} and
\S\ref{sec:composition}). Each environment is classified on four
axes (category, vector, salience, PII target) and a further four
factor axes (pressure cue, prompt-injection style, interaction
modality, multi-site flow), forming an 8-axis factorial taxonomy
mapped to published threat catalogues
\citep{mitre_attack_citation,owasp_llm_citation,enisa_citation}. 60
of the 91 environments are paired siblings that differ from a parent
on exactly one axis, supporting causal single-factor ablation. We
equip a Playwright-based agent
\citep{playwright,zhou2024webarenarealisticwebenvironment,Koh2024VisualWebArenaEM,10.5555/3692070.3694608,NEURIPS2023_5950bf29,he-etal-2024-webvoyager}
with a 40-field PII profile and evaluate four frontier model
families (GPT-5 mini, Claude Haiku 4.5, Gemini 3 Flash, Llama 4
Scout) across four mitigation conditions C0--C3 (baseline, generic
privacy nudge, phishing-aware checklist, pre-submission reflection).
The analysis plan is committed to a public git tag
(\texttt{prereg-v2-start}) before any reported data was collected.

\paragraph{Contributions and headline findings.} The paper makes
four contributions, each paired with the finding that backs it.

\begin{enumerate}[leftmargin=*,itemsep=3pt]
  \item \textbf{A detection--action gap under pre-submission
  reflection (primary).} Pooled across all four models in the
  reflection condition, agents whose reasoning an independent LLM
  judge confirms flagged the site as suspicious still leak critical
  PII at 35.9\%, versus 66.1\% otherwise---a 30.2\,pp gap
  ($q < 0.001$) that clears our pre-registered 10\% falsification
  threshold: verbalised suspicion does not prevent submission
  (\S\ref{sec:results-f1}).

  \item \textbf{Prompt-level mitigation is sharply model-dependent.}
  Baseline critical-tier leakage ranges from 54.5\% to 93.1\% across
  the four models (against 0\% on 10 benign twins), and escalating
  prompt-level mitigation reduces it substantially on three models
  but barely moves the fourth (Llama). Three model--condition cells
  individually cross the pre-registered $-$30\,pp threshold while the
  pooled effect ($-$23.3\,pp) does not: the heterogeneity, not a
  single pooled number, is the finding, and it motivates output-level
  defences over further system-prompt iteration
  (\S\ref{sec:results-mitigation}).

  \item \textbf{A characterisation of where leakage concentrates.}
  Across the 8-axis taxonomy no factor axis survives BH correction:
  paired-sibling toggles of salience, pressure, prompt-injection
  style, and interaction modality move pooled
  PLR\textsubscript{crit} by at most a few points. The
  methodological payoff is a sign disagreement on salience: a
  cross-cutting marginal ranks subtle attacks most dangerous, while
  the paired-sibling test (controlling for category base rates)
  finds subtle siblings leak \emph{less} than their parents,
  illustrating why the paired design is necessary
  (\S\ref{sec:results-axes}).

  \item \textbf{The \textsc{Scammer4U} benchmark}, released in full:
  91 attacker-controlled environments and 10 benign twins on an
  8-axis taxonomy mapped to published threat catalogues, with a
  pre-registered analysis plan and complete run logs, so follow-up
  work can re-run these ablations on new models without re-authoring
  the apparatus.
\end{enumerate}

